% !Mode:: "TeX:UTF-8"
% !Mode:: "TeX:UTF-8"
\newcommand{\thepage}{} % to silence a hyperref warning about pdfpagelabels
\RequirePackage[]{silence}
\WarningsOff[hyperref]
\PassOptionsToPackage{log-declarations=false}{xparse}
%\RequirePackage[l2tabu, orthodox]{nag}
\RequirePackage{luatex85} % to make \pdf...= commands work with LuaTeX
\documentclass[aspectratio=169,t]{beamer}
\usepackage{ifthen}
\usepackage[T1]{fontenc}
\usepackage{iftex}
\ifLuaTeX
 \usepackage{fontspec}
 \setmainfont[Ligatures=TeX]{Source Sans Pro}
\else
 \usepackage[utf8]{inputenc}
\fi
\usepackage{german}
\usepackage{newunicodechar}
\newunicodechar{€}{\texteuro}
\newunicodechar{ß}{\ss}
\renewcommand{\SS}{\iffontchar\font"1E9E \symbol{"1E9E}\else SS\fi}
\usepackage{ellipsis}
\usepackage{microtype}
\usepackage{listings}
\usepackage{sourcecodepro}
\lstset{language=Python, basicstyle=\ttfamily}
\usepackage{amsfonts}
\usepackage{amssymb}
\usepackage{amsthm}
\usepackage{fancyhdr}
\usepackage{graphicx}

\newtheorem{proposition}[theorem]{Proposition}

\DeclareMathOperator*{\argmax}{argmax}
\DeclareMathOperator{\supp}{supp}
\newcommand{\E}{\ensuremath{\mathbb{E}}}
\newcommand{\V}{\ensuremath{\mathbb{V}}}

\setbeamercolor{emph}{fg=purple}
\renewcommand<>{\emph}[1]{%
  {\usebeamercolor[fg]{emph}\only#2{\itshape}#1}%
}



\hypersetup{draft=false,colorlinks=false,pdfpagemode=UseOutlines,plainpages=false,
            breaklinks=true,bookmarksnumbered=true,pdfborder=0 0 0,
            pdftitle={Data Science SoSe 2026 HHU},
            pdfauthor={Konrad Voelkel},
            pdfsubject={data science}}

%\documentclass[aspectratio=169,t]{beamer}
\usetheme{hhubeamer} % put the hhu-beamer-theme in /usr/share/texmf/tex/latex/beamer/themes/

% -------------------- Code listings config --------------------
% Define custom colors for syntax highlighting
\definecolor{codegreen}{rgb}{0,0.6,0}
\definecolor{codegray}{rgb}{0.5,0.5,0.5}
\definecolor{codepurple}{rgb}{0.58,0,0.82}
\definecolor{backcolour}{rgb}{0.96,0.96,0.94} % Very light gray/beige
% Configure the listings package
\lstset{
    language=Python,
    backgroundcolor=\color{backcolour},   
    commentstyle=\color{codegreen}\itshape,
    keywordstyle=\color{blue}\bfseries,
    numberstyle=\tiny\color{codegray},
    stringstyle=\color{codepurple},
    basicstyle=\ttfamily\footnotesize, % Use a monospaced font
    breakatwhitespace=false,         
    breaklines=true,      
             % Enable automatic line breaking
    captionpos=b,                      % Caption at the bottom
    keepspaces=true,                 
    numbers=left,                      % Line numbers on the left
    numbersep=10pt,                    % Distance of line numbers from the code
    showspaces=false,                
    showstringspaces=false,            % Don't show weird underscores in strings
    showtabs=false,                  
    tabsize=4,
    frame=single,                      % Add a subtle border around the code
    rulecolor=\color{black!20},        % Border color
  }

\event[]{}
\author[Sommersemester 2026]{Dr. Konrad Völkel, Sommersemester 2026}

%\title{Data Science, VL 1}
%\date{16.04.2026}


%%% Local Variables:
%%% mode: latex
%%% TeX-master: t
%%% End:


\title{Data Science, VL 17}
\subtitle{Maximum Likelihood}
\date{23.06.2026}

\begin{document}

\begin{frame}
  \titlepage
\end{frame}

\section{Dimensionsreduktion mit PCA}

\begin{frame}[separator]
  \frametitle{Herleitung der PCA}
\end{frame}


\begin{frame}
\frametitle{PCA}
\alert{Principal Component Analysis} (PCA, Hauptkomponentenanalyse) ist ein \emph{unsupervised} Verfahren; \\
meist sind damit Schritte 1-2 oder auch 1-3 gemeint:

\begin{enumerate}[<+->]
\item (fit) Berechnung einer linearen invertierbaren Koordinatentransformation \\(einer OrthoNormalBasis) aus Daten
\item (transform) Anwenden der Transformation auf Daten
\item(Dimensionsreduktion) Projektion der Daten auf einen niedrigdimensionalen linearen Unterraum 
\end{enumerate}

\uncover<4->{
  Dabei werden die Daten so transformiert, dass die neuen Koordinatenachsen (sog. Hauptkomponenten/Hauptachsen) geordnet sind danach, wo die Daten die höhere Varianz aufweisen.
  }
\end{frame}

\begin{frame}{Beispiel}\framesubtitle{(1/2)}
  \vspace*{-2.3cm}

  \hspace{-3cm}
\includegraphics[height=1.1\textheight]{figures/case_a_dual_proj.png}
\end{frame}

\begin{frame}{Beispiel}\framesubtitle{(2/2)}
  \vspace*{-2.3cm}

  \hspace{-3cm}
\includegraphics[height=1.1\textheight]{figures/case_b_dual_proj.png}
\end{frame}

\begin{frame}
  \frametitle{Beispiel}
  \framesubtitle{Die vorigen Plots in Worte gefasst}
Wenn man Daten im dreidimensionalen Raum hat, bei denen aber die dritte Koordinate $z$
 stets $1$ ist, so ändert eine Projektion auf die $xy$-Ebene \alert{nichts} am Informationsgehalt der Daten.

\uncover<2->{
Wenn die Daten hingegen in der $z$-Koordinate geringfügig um die $1$ variieren, in den $x$- und $y$-Koordinaten aber deutlich stärker variieren, so ändert eine Projektion auf die 
$xy$-Ebene nur \alert{wenig} am Informationsgehalt der Daten.
}
\end{frame}


\begin{frame}
  \frametitle{PCA als Koordinatentransformation}
  \begin{itemize}[<+->]
   \item Eine Koordinatentransformation ist genau die Matrixmultiplikation mit einer invertierbaren Matrix.
\item  Die Spalten der Matrix sind die Bilder der Basisvektoren.
\item  Die PCA entspricht einer Orthonormalmatrix, d.h. die Spalten sind normiert und stehen senkrecht aufeinander.
\end{itemize}
\end{frame}

\begin{frame}
  \frametitle{Richtung der größten Varianz}
  \begin{itemize}[<+->]
  \item Ziel der PCA 1) \\Richtung (Vektor, Hauptachse) finden sodass in der Richtung die höchste Varianz vorliegt.
  \item Ziel der PCA 2) \\Für das Orthogonalkomplement der Hauptachse wiederum PCA.
  \item Ergebnis: \\Liste mit Hauptachsen (=Matrix), die senkrecht aufeinander stehen (=Orthogonalmatrix).
  \item Transformation auf die neue Basis mit dem \alert{Inversen} dieser Matrix.
  \end{itemize}
  \end{frame}


\begin{frame}
\frametitle{Beispiel}
Wenn wir ein Sample einer bivariaten Standardnormalverteilung betrachten, können wir einen Kreis um das 95\%-Quantil zeichnen.
Die Matrix $\begin{pmatrix} 2 & 0 \\ 0 & 1\end{pmatrix}$ ist eine Streckung auf der $x$-Achse, die aus dem hübschen Kreis eine langgezogene Ellipse macht.
Die Richtung der größten Varianz ist dann genau die $x$-Achse.

\uncover<2->{
Wenn eine beliebige invertierbare Matrix $M = \begin{pmatrix} a & b \\ c & d \end{pmatrix}$ auf ein Sample einer bivariaten Standardnormalverteilung angewendet wird, dann können wir die Richtungen mit größter Varianz bestimmen als die Richtungen, die von $M$ am stärksten gestreckt werden. Dass wäre ein Vektor $v$ mit $Mv = \lambda v$ für das größte solche $\lambda$.
}
\end{frame}


\begin{frame}
\frametitle{Beobachtung}
Wenn wir die Kovarianzmatrix einer bivariaten Standardnormalverteilung betrachten, also $\begin{pmatrix} \mathbb{V}(X) & Cov(X,Y) \\ Cov(Y,X) & \mathbb{V}(Y) \end{pmatrix} = \begin{pmatrix} 1 & 0 \\ 0 & 1 \end{pmatrix}$
und dann, wie sich die Kovarianzmatrix eines Samples ändert, wenn wir die Koordinatentransformation $\begin{pmatrix} 2 & 0 \\ 0 & 1\end{pmatrix}$ durchführen,
beobachten wir eine Quadrierung: beim Berechnen der Sample-Kovarianz $Cov(x,y) = \frac{1}{n} \sum_{i=1}^n x_iy_i$ geht die Skalierung mit $2$ einfach ein, bei $\mathbb{V}(x) = Cov(x,x)$ aber quadratisch.

\uncover<2->{
Wenn wir also die Wurzel aus der neuen Kovarianzmatrix berechnen, erhalten wir die Koordinatentransformation zurück.
}
\end{frame}


\begin{frame}
\frametitle{Matrixwurzel}
Kovarianzmatrizen sind positiv definit und symmetrisch.

\uncover<2->{
  Zu jeder positiv definiten symmetrischen Matrix $C$\\
  existiert eine invertierbare Matrix $\Sigma$ mit $\Sigma^2 = C$.
  }
  \uncover<3->{

\vspace{1cm}
    
Eine quadratische Matrix $C$ heisst \\ symmetrisch, wenn $C=C^T$ und \\ positiv definit, wenn alle Eigenwerte positiv sind.

\vspace{1cm}

Äquivalent: $C$ positiv definit wenn für alle $v \neq 0$ gilt dass $v^TCv > 0$ ist.
    }
\end{frame}


\begin{frame}
\frametitle{Zutaten zusammensetzen}
Wir stellen uns jetzt vor, dass die Daten $X$ aus einer Standardnormalverteilung $Y$ durch eine Koordinatentransformation $A$ hervorgehen.

\begin{itemize}[<+->]
  \item Können wir aus $X = AY$ etwa $A$ oder $Y$ zurückgewinnen? \\ Was sind die Hauptkomponenten?
  \item  Die empirische Kovarianzmatrix schätzt die Kovarianzmatrix $C$ ein,\\
  deren Wurzel $A$ ist, also $A^2 = C$.
\item Wir berechnen den größten Eigenwert von $A$ und einen zugehörigen Eigenvektor $v$.
\item Das ist die Richtung der größten Varianz in den Daten! 
\item Es ist auch ein Eigenvektor zum größten Eigenwert von $A^2$.
\end{itemize}
\end{frame}


\begin{frame}
\frametitle{Letzter Schritt}
Wenn die Daten nicht normalverteilt sind (oder nur $A$ unbekannt ist),
können wir das gleiche machen:

\uncover<2->{
Einen Eigenvektor zum größten Eigenwert von der empirischen Kovarianzmatrix $C$ bestimmen.
}

\uncover<3->{
Schön: Das ist \alert{immer} die Richtung der größten Varianz in den Daten.
}
\end{frame}


\begin{frame}
\frametitle{Anderer Zugang: Optimierung}
Die Richtung größter Varianz lässt sich auch finden als die Gerade,
welche die Summe der quadrierten Abstände der Daten zur Gerade minimiert.

\uncover<2->{
  Das ist äquivalent zum Minimieren des MSE.
  Erinnerung: das konstante Modell, welches den MSE minimiert, ist \texttt{mean}.

  Die erste Hauptkomponente ist also das beste lineare Modell aus der MSE-Perspektive.
}

\uncover<3->{
\vspace{1cm}
  
(im Skript ist eine Animation dazu)
}
\end{frame}


\begin{frame}
\frametitle{Ab jetzt ins Notebook!}
\end{frame}


\begin{frame}[separator]
  \frametitle{Vielen Dank für Ihre Aufmerksamkeit!}
\end{frame}

\end{document}

%%% Local Variables:
%%% mode: latex
%%% TeX-master: t
%%% End:
